\section{引言}
\sourcepage{301}
我们已经给每个拓扑空间 $X$ 指定了一列阿贝尔群，称为同调群。现在再给它指定另一列阿贝尔群，称为\textbf{上同调群}（cohomology groups）。上同调的定义比同调晚得多，原因不难理解：它在几何上远不如同调自然，其起源在代数而非几何。在稍后将精确说明的某种代数意义下，它与同调互为“对偶”。过去拓扑学家曾把这些群元素的代表称为“伪闭链”等，似乎对它们能否成为正当研究对象抱有疑虑。后来逐渐明白，上同调既有理论意义，也有实际用途；它比同调更细致，也更深刻、更有力。本章将研究它。

流形的对偶定理、拓扑与微分几何之间的联系（de Rham 定理）、拓扑与分析之间的联系（层系数上同调），都用上同调表述。甚至空间的同胚分类、映射的同伦分类等纯拓扑问题，上同调也能提供丰富信息。以后将回到其中一些问题。

本章假定读者熟悉第 3.7 节的范畴与函子语言。
